\emph{(i)} For $n \ge 2$ the relation for the pair $(1, 2)$ and the relation for the pair $(2, 1)$ give $e_{1}e_{2} = c\,e_{2}e_{1} = c^{2}e_{1}e_{2}$, so $c^{2} = 1$ by cancellation in $G$ (the same argument as the sign from symmetry, clause (i) of Theorem \claimref{synthesis.blade-square-signs-f2-quadratic-form}, read in the group). For $n = 1$ there is no pair. In the cyclic group of order $2m$ generated by $e_{1}$, with $m \ge 2$ and $c = e_{1}^{2}$, the element $c$ is central and $\ne 1$, $G/\langle c\rangle = \{\langle c\rangle, e_{1}\langle c\rangle\}$ has order $2$ with basis the image of $e_{1}$, and $c$ has order $m$, so $c^{2} = 1$ exactly when $m = 2$. Existence: $f(S,T) = \sum_{i > j}x_{i}y_{j} + \sum_{i}s_{i}x_{i}y_{i}$, with $x, y \in \mathbb{F}_{2}^{n}$ the indicator vectors of $S, T$, is $\mathbb{F}_{2}$-bilinear. So on $\{0, 1\} \times \mathbb{F}_{2}^{n}$ the product $(a, S)(b, T) = (a + b + f(S,T), S \triangle T)$ (first entry modulo $2$) is associative, both triple products being $(a + b + d + f(S,T) + f(S,U) + f(T,U), S \triangle T \triangle U)$ for a third factor $(d, U)$; $(0, \emptyset)$ is the identity, as $f$ vanishes when either argument is $\emptyset$; and $(a + f(S,S), S)$ is the inverse of $(a, S)$. Put $c = (1, \emptyset)$, $e_{i} = (0, \{i\})$. Then $c$ is central with $c^{2} = 1 \ne c$; $e_{i}^{2} = (f(\{i\}, \{i\}), \emptyset) = c^{s_{i}}$; for $i \ne j$ the products $e_{i}e_{j}$ and $e_{j}e_{i}$ are $(f(\{i\},\{j\}), \{i, j\})$ and $(f(\{j\},\{i\}), \{i, j\})$, whose first entries add to $1$, so $e_{i}e_{j} = c\,e_{j}e_{i}$; and since $f(S, \{j\}) = 0$ when $j$ exceeds every element of $S$, the increasing product $e_{S}$ is $(0, S)$, so $(a, S) \mapsto S$ identifies the quotient by $\langle c\rangle$ with $\mathbb{F}_{2}^{n}$, the $e_{i}$ mapping to the basis. This is a group of signed directions with the given signs; uniqueness up to an isomorphism fixing $c$ and the $e_{i}$ follows from the presentation below. Granting $c^{2} = 1$ and $c \ne 1$, $\langle c\rangle = \{1, c\}$; each coset of $\langle c\rangle$ is $e_{S}\langle c\rangle$ for exactly one $S$, so each element is $c^{a}e_{S}$ for exactly one $a \in \{0, 1\}$ and one $S$, and $|G| = 2^{n+1}$. The image of $e_{i}^{2}$ in $\mathbb{F}_{2}^{n}$ is $\{i\} \triangle \{i\} = \emptyset$, so $e_{i}^{2} \in \{1, c\}$. Let $\Gamma$ be the group presented by the generators and relations of the statement. In $\Gamma$, $c^{-1} = c$ and $e_{i}^{-1} = c^{s_{i}}e_{i}$, so every element is a word in $c, e_{1}, \dots, e_{n}$ without inverses, and the rewriting of (ii), which uses only the relations, brings it to the form $c^{a}e_{S}$; hence $|\Gamma| \le 2^{n+1}$. The elements $e_{i}$ and $c$ of $G$ satisfy the relations and generate $G$, so the homomorphism $\Gamma \to G$ sending generators to generators is onto, and since $|G| = 2^{n+1} \ge |\Gamma|$ it is an isomorphism.

\emph{(ii)} In the word $e_{S}e_{T} = e_{i_{1}}\cdots e_{i_{k}}e_{j_{1}}\cdots e_{j_{l}}$ ($i_{1} < \dots < i_{k}$, $j_{1} < \dots < j_{l}$) move $e_{j_{1}}, e_{j_{2}}, \dots$ in turn to the left until they reach their place in increasing order. The letter $e_{j}$ passes exactly the letters $e_{i}$ of $e_{S}$ with $i > j$, each passage costing a factor $c$ (the relation, with $c$ central), and if $j \in S$ it then meets $e_{j}$ and the two merge into $e_{j}^{2} = c^{s_{j}}$; it never meets a letter of $e_{T}$ already placed, those being smaller. Collecting the central factors gives $e_{S}e_{T} = c^{f(S,T)}e_{S \triangle T}$, the sign rule of the blade product of an orthogonal basis, with $c$ in place of $-1$. Since $G/\langle c\rangle$ is elementary abelian, every square $g^{2}$ and every commutator $ghg^{-1}h^{-1}$ maps to $0$ and lies in $\langle c\rangle = \{1, c\}$, which is central. For $g = c^{a}e_{S}$, $g^{2} = c^{2a}e_{S}^{2} = c^{f(S,S)}$, and $f(S,S) = \binom{|S|}{2} + |S \cap N|$, the pairs $(i, j) \in S \times S$ with $i > j$ being the $\binom{|S|}{2}$ two-element subsets of $S$. Also $e_{T}e_{S} = c^{f(T,S)}e_{S \triangle T}$, so, as $c^{-1} = c$, $e_{S}e_{T} = c^{f(S,T) + f(T,S)}e_{T}e_{S}$, and $f(S,T) + f(T,S)$ counts the ordered pairs $(i, j) \in S \times T$ with $i \ne j$, plus $2|S \cap T \cap N|$: it is $|S|\,|T| - |S \cap T| \equiv |S|\,|T| + |S \cap T| = b(S,T) \bmod 2$. The central factors $c^{a}$ do not change squares or commutators. From $hg = c^{b(S,T)}gh$, $(gh)^{2} = g(hg)h = c^{b(S,T)}g^{2}h^{2}$, so $q(S \triangle T) = q(S) + q(T) + b(S,T)$: $q$ is a quadratic form over $\mathbb{F}_{2}$ with polar form $b = f + f^{\top}$, which is bilinear because $f$ is; these are the same forms as in clause (ii) of Theorem \claimref{synthesis.blade-square-signs-f2-quadratic-form} (there with the negative directions numbered last; the formulas do not depend on the numbering). For $n \ge 2$, $e_{1}e_{2}e_{1}^{-1}e_{2}^{-1} = c$ is a commutator. It is also a square: if $s_{i} = s_{j}$ for some $i \ne j$, then $(e_{i}e_{j})^{2} = c^{q(\{i, j\})} = c^{1 + s_{i} + s_{j}} = c$; otherwise $n = 2$ and $\{s_{1}, s_{2}\} = \{0, 1\}$, and $e_{i}^{2} = c$ for the $i$ with $s_{i} = 1$.

\emph{(iii)} An element $c^{a}e_{S}$ is central exactly when it commutes with every $e_{i}$, since the $e_{i}$ and $c$ generate $G$, that is, when $b(S, \{i\}) = 0$ for every $i$: when $S$ lies in the radical of $b$. Here $b(S, \{i\}) = |S| + [i \in S]$. If $S$ is neither empty nor $\{1, \dots, n\}$, taking $i \in S$ and $i' \notin S$ gives $|S| \equiv 1$ and $|S| \equiv 0$, which is impossible; $\emptyset$ is always in the radical, and $\{1, \dots, n\}$ is exactly when $n + 1 \equiv 0$, that is $n$ odd. So the radical is $\{\emptyset\}$ for $n$ even and $\{\emptyset, \{1, \dots, n\}\}$ for $n$ odd (as in clause (iii) of Theorem \claimref{synthesis.blade-square-signs-f2-quadratic-form}), which gives the centre. For $n$ even, the centre, the commutators and the squares are all $\{1, c\}$ by (ii) ($n \ge 2$), so the subgroup they generate is $\{1, c\}$ as well, and $G$ is extraspecial of order $2^{1+n}$ (for a finite $2$-group the subgroup generated by squares and commutators is the Frattini subgroup). Central elements of order $2$: $c$, and for $n$ odd $c^{a}e_{1 \cdots n}$ exactly when $(c^{a}e_{1 \cdots n})^{2} = c^{q(\{1, \dots, n\})} = 1$. Here $q(\{1, \dots, n\}) = \binom{n}{2} + q$, and for $n = p + q$ odd, $\binom{n}{2} + q$ is even exactly when $p - q \equiv 1 \bmod 4$: writing $p - q = n - 2q$, both conditions are unchanged when $q$ grows by $2$ and both flip when $q$ grows by $1$ (since $n - 2q$ then drops by $2$), and for $q = 0$ the first reads $n \equiv 1 \bmod 4$, which is the second.

\emph{(iv)} Let $\Phi : G \to G'$ be an isomorphism with $\Phi(c) = c'$. It induces a group isomorphism $\phi : G/\langle c\rangle \to G'/\langle c'\rangle$, that is, a linear bijection of $\mathbb{F}_{2}^{n}$, and $\Phi(g)^{2} = \Phi(g^{2})$ reads $c'^{\,q'(\phi(S))} = c'^{\,q(S)}$, so $q' \circ \phi = q$. Conversely, let $q' \circ \phi = q$. Then also $b' \circ (\phi \times \phi) = b$, $b$ being the polar form of $q$. Choose $g_{i} \in G'$ over $\phi(\{i\})$. By (ii) for $G'$, $g_{i}^{2} = c'^{\,q'(\phi(\{i\}))} = c'^{\,s_{i}}$ and $g_{i}g_{j} = c'^{\,b(\{i\}, \{j\})}g_{j}g_{i} = c'\,g_{j}g_{i}$ for $i \ne j$, and $c'$ is central with $c'^{2} = 1$. By the presentation of (i) there is a homomorphism $\Phi : G \to G'$ with $\Phi(e_{i}) = g_{i}$ and $\Phi(c) = c'$; its image contains $c'$ and elements over the basis $\phi(\{1\}), \dots, \phi(\{n\})$, hence every coset of $\langle c'\rangle$ and all of $G'$; so $\Phi$ is onto, and an isomorphism because $|G| = |G'| = 2^{n+1}$. For $n \ge 2$, $c$ is the only commutator other than $1$ in $G$ and $c'$ in $G'$, and isomorphisms carry commutators to commutators.

\emph{(v)} $\Z[G]$ is free on the $2^{n+1}$ elements $c^{a}e_{S}$. As $c$ is central, the two-sided ideal generated by $c + 1$ is $(c + 1)\Z[G]$, which is spanned over $\Z$ by the $(c + 1)c^{a}e_{S}$; since $c(c + 1) = c + c^{2} = c + 1$, these are the $(1 + c)e_{S}$. So $\Z[G] = \bigoplus_{S}(\Z e_{S} \oplus \Z\,c\,e_{S})$ and $(c + 1) = \bigoplus_{S}\Z(e_{S} + c\,e_{S})$; in each summand $e_{S}$ and $e_{S} + c\,e_{S}$ form a $\Z$-basis (the change of basis from $e_{S}, c\,e_{S}$ has determinant $1$), so $A = \bigoplus_{S}\Z e_{S}$, free with basis the images of the $e_{S}$, with $c\,e_{S} \mapsto -e_{S}$. The product follows from (ii).

\emph{(vi)} $v^{2} = \sum_{i}a_{i}^{2}e_{i}^{2} + \sum_{i < j}a_{i}a_{j}(e_{i}e_{j} + e_{j}e_{i})$, and $e_{i}e_{j} + e_{j}e_{i} = (c + 1)e_{j}e_{i}$ maps to $0$ in $A$, while $e_{i}^{2} = c^{s_{i}}$ maps to $(-1)^{s_{i}}$. So $v^{2} = \sum_{i}(-1)^{s_{i}}a_{i}^{2} = Q(v)$; no element $\tfrac12$ is used.

\emph{(vii)} Tensoring the exact sequence $(c + 1) \to \Z[G] \to A \to 0$ with $R$ gives $A_{R} = R[G]/(c + 1)$, free over $R$ on the $1 \otimes e_{S}$, which are products of the $\iota(\epsilon_{i}) = e_{i}$ ($\epsilon_{i}$ the standard basis of $R^{n}$). Uniqueness: $A_{R}$ is generated as an $R$-algebra by $\iota(R^{n})$. Existence: put $u_{i} = f(\epsilon_{i})$. Then $u_{i}^{2} = Q(\epsilon_{i})\,1 = (-1)^{s_{i}}$, and for $i \ne j$, $(u_{i} + u_{j})^{2} = Q(\epsilon_{i} + \epsilon_{j}) = (-1)^{s_{i}} + (-1)^{s_{j}}$ gives $u_{i}u_{j} + u_{j}u_{i} = 0$. Each $u_{i}$ is a unit, with inverse $(-1)^{s_{i}}u_{i}$, and in the unit group of $B$ the elements $u_{1}, \dots, u_{n}$ and $-1$ satisfy the relations of the presentation in (i): $-1$ is central with square $1$, $u_{i}^{2} = (-1)^{s_{i}}$, $u_{i}u_{j} = (-1)u_{j}u_{i}$. So there is a group homomorphism $\psi : G \to B^{\times}$ with $\psi(e_{i}) = u_{i}$ and $\psi(c) = -1$; its $R$-linear extension $R[G] \to B$ is a unital algebra homomorphism that kills $c + 1$, hence factors through $A_{R}$ as $F$, and $F(\iota(x)) = \sum_{i}x_{i}u_{i} = f(x)$. The tensor algebra of $R^{n}$ modulo the ideal generated by the $x \otimes x - Q(x)$ has the same universal property (that of the tensor algebra, passed to the quotient), so the two are isomorphic by the unique isomorphism compatible with the maps from $R^{n}$; $\iota$ is injective because the $e_{i}$ are part of an $R$-basis. For $R = \R$ the universal property is that of Definition \claimref{obj.std-algebra.clifford-algebra} for $\R^{n}$ with the form $Q$ of signature $(p, q)$. If $2 = 0$ in $R$ then $-1 = 1$, so $e_{S}e_{T} = e_{S \triangle T}$, which is the multiplication of the group algebra of $\mathbb{F}_{2}^{n}$.

\emph{(viii)} (a) $\Leftrightarrow$ (b) is (iv). (b) $\Rightarrow$ (c): an isomorphism $G \to G'$ carrying $c$ to $c'$ extends to $R[G] \cong R[G']$ carrying $c + 1$ to $c' + 1$. (c) $\Rightarrow$ (d) with $R = \Z$. (d) $\Rightarrow$ (e): let $R \to K$ be a unital homomorphism to an ordered field; then $K \otimes_{R} A_{R} = A_{K}$, so an isomorphism over $R$ gives one over $K$. On $A_{K}$ put $\tau(x, y) = 2^{-n}\Tr(L_{xy})$, the trace of left multiplication on the $K$-vector space $A_{K}$. An algebra isomorphism $\Phi$ satisfies $L_{\Phi(z)} = \Phi L_{z}\Phi^{-1}$, so $\tau$ is carried to $\tau'$. Since $e_{S}e_{T} = \pm e_{S \triangle T}$, $L_{e_{S}}$ maps each basis line to a different one unless $S = \emptyset$, so $\Tr(L_{e_{S}}) = 2^{n}$ for $S = \emptyset$ and $0$ otherwise; hence $\tau(x, y) = \langle xy\rangle_{0}$, $\tau(e_{S}, e_{T}) = 0$ for $S \ne T$, and $\tau(e_{S}, e_{S}) = (-1)^{q(S)}$. Let $P$ be the number of $S$ with $q(S) = 0$ and $U_{\pm}$ the spans of the blades with $e_{S}^{2} = \pm 1$. On $U_{+}$, $\tau(x, x) = \sum x_{S}^{2} > 0$ for $x \ne 0$, and on $U_{-}$, $\tau(x, x) = -\sum x_{S}^{2} < 0$ for $x \ne 0$; a subspace $W$ with $\dim W > P$ meets $U_{-}$ in a nonzero vector, as $\dim W + \dim U_{-} > 2^{n}$. So $P$ is the largest dimension of a subspace on which $\tau$ is positive definite, an isomorphism invariant, and with it the Gauss sum $\sum_{S}(-1)^{q(S)} = 2P - 2^{n}$. (e) $\Rightarrow$ (a): all the forms $q_{(p,q)}$ with the same $n$ have the same polar form $b$, whose radical is $\{0\}$ for $n$ even and $\{0, z\}$, $z = \{1, \dots, n\}$, for $n$ odd. Let $W$ be a complement of the radical ($W = \mathbb{F}_{2}^{n}$ for $n$ even); $b$ is alternating ($b(x, x) = q(0) = 0$) and nondegenerate on $W$, so $W$ has a basis $u_{1}, w_{1}, \dots, u_{m}, w_{m}$ with $b(u_{k}, w_{k}) = 1$ and $b = 0$ on all other pairs of basis vectors (choose $u \ne 0$, then $w$ with $b(u, w) = 1$, and continue in the $b$-orthogonal complement of $\mathrm{span}\{u, w\}$, which is complementary to it and on which $b$ is nondegenerate). On each plane $\mathrm{span}\{u_{k}, w_{k}\}$ the three nonzero values of $q$ add to $1$, since $q(u + w) = q(u) + q(w) + 1$: either two of them are $0$, and the two vectors with $q = 0$ form a new basis of the plane with $b = 1$ (type $H$, Gauss sum $2$), or all three are $1$ (type $P$, Gauss sum $-2$). Gauss sums multiply over $b$-orthogonal sums, since $q(x + y) = q(x) + q(y)$ when $b(x, y) = 0$. Two planes of type $P$ give one of type $H$ and one more: in $\mathrm{span}\{u_{1}, w_{1}, u_{2}, w_{2}\}$, $u = u_{1} + u_{2}$ and $w = w_{1} + u_{2}$ have $q(u) = q(w) = 0$ and $b(u, w) = 1$, and the orthogonal complement of $\mathrm{span}\{u, w\}$ there is a plane of Gauss sum $(-2)(-2)/2 = 2$, so of type $H$. Hence $q|_{W}$ is equivalent to $H^{m}$, Gauss sum $2^{m}$, or to $H^{m-1} \perp P$, Gauss sum $-2^{m}$. For $n$ even this is $q$ itself. For $n$ odd, $q(x + az) = q(x) + a\,q(z)$ for $x \in W$, as $z$ is in the radical. If $q(z) = 1$, the pairing $x \leftrightarrow x + z$ shows that the Gauss sum is $0$, and a plane of type $P$ in $W$ becomes one of type $H$ when $u$ is replaced by $u + z$ ($q(u + z) = 0$, $b$ unchanged), so $q$ is equivalent to $H^{m} \perp \langle z\rangle$ with $q(z) = 1$; if $q(z) = 0$, $q$ is equivalent to $q|_{W} \perp \langle z\rangle$ with $q(z) = 0$, with Gauss sum $\pm 2^{m+1}$, the sign that of $q|_{W}$. In each parity the models have distinct Gauss sums, so two forms with equal Gauss sums are equivalent to the same model, hence to each other. The values: in $\Cl_{\Z}(0,1) = \Z \oplus \Z j$, $j^{2} = -1$, a commutative ring, $(-1)^{\binom{k}{2}} = \langle (1 - j)j^{k}\rangle_{0}$ (both sides are $1, 1, -1, -1$ for $k = 0, 1, 2, 3$ and have period $4$). Grouping the $S$ by $a = |S \setminus N|$ and $b' = |S \cap N|$, $(-1)^{q(S)} = (-1)^{\binom{a + b'}{2} + b'} = \langle (1 - j)j^{a}(-j)^{b'}\rangle_{0}$, so $\sum_{S}(-1)^{q(S)} = \langle (1 - j)(1 + j)^{p}(1 - j)^{q}\rangle_{0}$. With $(1 + j)(1 - j) = 2$, $(1 + j)^{2} = 2j$ and $(1 + j)^{8} = 16$: for $r = p - q = 8m + s \ge 0$ ($0 \le s \le 7$) the product is $2^{q}16^{m}(1 - j)(1 + j)^{s}$, and $\langle (1 - j)(1 + j)^{s}\rangle_{0} = 1, 2, 2, 0, -4, -8, -8, 0$ for $s = 0, \dots, 7$; for $r < 0$, $|r| = 8m + s'$, the product is $2^{p}(1 - j)^{|r| + 1}$, whose scalar part equals that of $2^{p}(1 + j)^{|r| + 1}$ (the automorphism $j \mapsto -j$), and $\langle (1 + j)^{t}\rangle_{0} = 1, 1, 0, -2, -4, -4, 0, 8, 16$ for $t = 0, \dots, 8$. With $n = r + 2q$ (respectively $|r| + 2p$) this is $2^{n/2}$ for $r \equiv 0, 2$, $-2^{n/2}$ for $r \equiv 4, 6$, $2^{(n+1)/2}$ for $r \equiv 1$, $-2^{(n+1)/2}$ for $r \equiv 5$ and $0$ for $r \equiv 3, 7 \pmod 8$, the values of clause (iv) of Theorem \claimref{synthesis.blade-square-signs-f2-quadratic-form}; for fixed $n$ they are distinct on the five classes, which are therefore the classes of (viii).

The check builds the cocycle group for every sign vector with $n \le 4$ and verifies associativity, inverses and the relations; it builds $G$ from its relations by rewriting words for every sign vector with $n \le 4$ and for the standard signatures with $n = 5$, and verifies the normal form, associativity ($n \le 3$), agreement with the product of signed blades in GA with $c \mapsto -1$, the closed forms of $q$ and $b$, and that the squares are $\{1, c\}$; the Heisenberg group modulo $k$ ($2 \le k \le 7$) reads its relation both ways only for $k = 2$, and the cyclic groups of order $2m$ ($2 \le m \le 7$) show the independence at $n = 1$; the centre, the central involutions and the extraspecial property for all $20$ signatures with $n \le 5$; the ideal, a unimodular basis change, integral structure constants equal to GA's and $v^{2} = Q(v)$ for all integer vectors with entries in $[-2, 2]$, $n \le 4$, and rational vectors; the universal property for $27$ maps into $\Cl(p{+}1, q{+}1)$ (twisted generators), rotor-conjugated frames and the grade involution; the reduction modulo $2$; and, for all pairs of sign vectors with $n \le 3$ and of standard signatures with $4 \le n \le 6$, that an explicit $\mathbb{F}_{2}$-isometry exists exactly when the Gauss sums agree, lifted for $n \le 4$ to anticommuting blades of the right squares whose products are $2^{n}$ distinct signed blades (an isomorphism over $\Z$), with the traces $\Tr(L_{e_{S}})$ and the diagonal of the trace form; and the radical and the scalar part of $(1 - j)(1 + j)^{p}(1 - j)^{q}$ in $\Cl_{\Z}(0,1)$, equal to the Gauss sum and to the tabulated value, for all signatures with $n \le 8$.
